\subsection{Signal Preprocessing and Network Training}
\label{sec:experiment-training}

The signal preprocessing pipeline applied to all acquired data follows the protocol described in Section~\ref{sec:preprocessing}. Raw laser ultrasonic signals acquired with limited shot averaging ($N = 16$) are processed through an FIR bandpass filter (100\,kHz--2\,MHz) to retain the Lamb wave frequency band, followed by a rectangular time-gate selecting a $30\,\mu\text{s}$ window centered on the predicted $S_0$ mode arrival. Per-shot $z$-score normalization is applied to the windowed segment to standardize amplitudes across acquisition sessions. All preprocessed signals are resampled to a fixed length $L = 1024$ samples to match the input dimensionality of PMDF-Net. The STFT parameters for the time-frequency branch (Hanning window $W = 256$, hop size $H = 64$, 75\% overlap) are held constant during both training and inference.

\textbf{Paired dataset construction.} Training data are assembled in paired format following the protocol in Section~\ref{sec:dataset}. Input signals are formed from $N = 16$ averaged laser ultrasonic acquisitions, while target signals are formed from $M = 256$ averaged acquisitions at the same spatial location. The $10\log_{10}(M/N) \approx 12$\,dB SNR difference between input and target establishes the denoising target for the network to learn.

\textbf{Data augmentation.} During training, each paired sample undergoes randomized per-sample augmentation to improve generalization: random time-shift by $\pm 12.5\%$ of signal length, additive Gaussian noise injection at SNR levels uniformly sampled from $[-6, +6]$\,dB relative to the clean signal, and amplitude scaling uniformly sampled from $[0.8, 1.2]$. Augmentation is applied independently to the limited-averaging input and the high-averaging target to preserve pairing integrity. No augmentation is applied during validation or testing.

\textbf{Training configuration.} PMDF-Net is trained using the AdamW optimizer~\cite{Loshchilov2019} with an initial learning rate of $1 \times 10^{-3}$, decoupled weight decay of $1 \times 10^{-4}$, and mini-batch size of $32$. The learning rate follows a cosine annealing schedule from $1 \times 10^{-3}$ to $1 \times 10^{-4}$ over the full training run. Dropout with probability $p = 0.2$ is applied at the fusion layers. Training terminates when the validation reconstruction loss shows no improvement for $15$ consecutive epochs, restoring the best-validated model weights. The model contains approximately $4.2 \times 10^6$ trainable parameters. On a single NVIDIA RTX 4090 (24\,GB memory), a full training run of $180$ epochs completes in approximately $4.5$ hours, yielding an average throughput of roughly $2.3 \times 10^4$ training steps per epoch.